\documentclass[UTF8]{ctexart}
\usepackage{xeCJK}
\usepackage{geometry}
\geometry{left=2cm,right=2cm,top=2.5cm,bottom=2.5cm}
\setlength{\headheight}{12.7pt}

\usepackage{titlesec}
\titleformat{\section}
  {\normalfont\large\bfseries}
  {\thesection}
  {1em}
  {}
\titlespacing*{\section}
  {0pt}
  {3.5ex plus 1ex minus .2ex}
  {2.3ex plus .2ex}

\usepackage{amsmath}
\usepackage{amssymb}
\usepackage{amsthm}
\usepackage{enumitem}
\setlist[enumerate]{itemsep=0pt,topsep=0pt,parsep=0pt,leftmargin=0pt,itemindent=2.5em}
\setlist[itemize]{itemsep=0pt,topsep=0pt,parsep=0pt,leftmargin=0pt,itemindent=2.5em}
\usepackage{fancyhdr}
\pagestyle{fancy}
\fancyhf{}
\usepackage{graphicx}
\usepackage{hyperref}
\usepackage{indentfirst}
\usepackage{lmodern}


\begin{document}
\setlength{\abovedisplayskip}{1pt}
\setlength{\belowdisplayskip}{1pt}

\section{偏序关系}

\hypertarget{sec:定义1.1}{}
\noindent\textbf{定义1.1 (偏序关系)}

给定集合$a$和$a$上的二元关系$\leqslant$. 如果以下三条成立:
\begin{enumerate}
    \item 自反性: $\forall x\in a:\,x\leqslant x$,
    \item 传递性: $\forall x,y,z\in a:\,
    x\leqslant y\land y\leqslant z\to x\leqslant z$,
    \item 反对称性: $\forall x,y\in a:\,x\leqslant y\land y\leqslant x\to x=y$,
\end{enumerate}
那么称$\leqslant$是$a$上的一个\textbf{偏序关系}.
也称$(a,\leqslant)$是一个\textbf{偏序结构}.

\vspace{10pt}

给定偏序结构$(a,\leqslant)$, 它衍生的$<$满足以下两条, 由此称之为\textbf{严格偏序}:
\begin{enumerate}
    \item 反自反性: $\forall x\in a:\,x\nless x$,
    \item 传递性: $\forall x,y,z\in a:\,x<y\land y<z\to x<z$.
\end{enumerate}
另外, 此时$\geqslant$也是偏序, $>$也是严格偏序.

\vspace{10pt}

\hypertarget{sec:例1.1}{}
\noindent\textbf{例1.1 (包含偏序)}

给定集合$a$, 包含关系$\{(x,y)\in a\times a\mid x\subset y\}$是$a$上的偏序关系,
称为\textbf{包含偏序}, 可以直接记作$\subset$.

\vspace{10pt}

从\href{01 预序关系#nameddest=sec:定义1.1}{预序关系}可以导出一个偏序关系.

\vspace{10pt}

\hypertarget{sec:定理1.1}{}
\noindent\textbf{定理1.1 (导出偏序)}

给定预序结构$(a,\rho)$, 由\href{01 预序关系#nameddest=sec:定理1.1}{引理}有等价关系
$\sim\,=\{(x,y)\in a\times a\mid x\rho y\land y\rho x\}$, 那么
\[
    \leqslant\,=\left\{(X,Y)\in
    \,^{\displaystyle a}\!/_{\displaystyle\sim}\,\times
    \,^{\displaystyle a}\!/_{\displaystyle\sim}\,\middle\vert\,
    \exists x\in X\,\exists y\in Y:x\rho y\right\}
\]
是$^{\displaystyle a}\!/_{\displaystyle\sim}$上的偏序关系.
此时对任意的$x,y\in a$有$x\rho y\Leftrightarrow[x]\leqslant[y]$.

\noindent{\kaishu 证明.}

\begin{enumerate}
    \item 对任意的$X\in\,^{\displaystyle a}\!/_{\displaystyle\sim}$,
    取$x\in X$则有$x\rho x$, 即$X\leqslant X$.
    \item 对任意的$X,Y,Z\in\,^{\displaystyle a}\!/_{\displaystyle\sim}$,
    如果$X\leqslant Y$且$Y\leqslant Z$, 那么存在$x\in X,\ y,y'\in Y,\ z\in Z$使得
    \[
        x\rho y\quad\land\quad y'\rho z,
    \]
    再由$y\sim y'$得$y\rho y'$, 从而$x\rho z$, 即$X\leqslant Z$.
    \item 对任意的$X,Y\in\,^{\displaystyle a}\!/_{\displaystyle\sim}$,
    如果$X\leqslant Y$且$Y\leqslant X$, 那么存在$x,x'\in X,\ y,y'\in Y$使得
    \[
        x\rho y\quad\land\quad y'\rho x',
    \]
    再由$x\sim x',\,y\sim y'$得$y\rho y',\,x'\rho x$, 从而$y\rho x$.
    依定义有$x\sim y$, 于是$X=Y$.
\end{enumerate}

进一步, 如果$[x]\leqslant[y]$, 那么存在$x'\in[x],\ y'\in[y]$使得$x'\rho y'$,
显然$x\rho y$. {\hfill $\qedsymbol$}






\newpage

\hypertarget{sec:定理1.2}{}
\noindent\textbf{定理1.2 (偏序集表示定理)}

任给偏序结构$(a,\leqslant)$, 存在集合$b$使得$(a,\leqslant)\cong(b,\subset)$.

\noindent{\kaishu 证明.}

对任意的$x\in a$定义
\[
    f(x)=\{t\in a\mid t\leqslant x\},
\]
令$b=\{f(x)\mid x\in a\}$, 则有满射$f:a\to b$.

接下来, 对任意的$x,y\in a$, 首先如果$x\leqslant y$, 那么对任意的$t\in f(x)$有
\[
    t\in f(x)\implies t\leqslant x\implies
    t\leqslant y\implies t\in f(y),
\]
即$f(x)\subset f(y)$. 其次如果$f(x)\subset f(y)$, 那么
\[
    x\in f(x)\implies x\in f(y)\implies x\leqslant y,
\]
综上即得$x\leqslant y\leftrightarrow f(x)\subset f(y)$.

最后, 如果$f(x)=f(y)$, 那么同时有$x\leqslant y$和$y\leqslant x$,
即$x=y$. 所以$f$是单射.

综上即得$f:(a,\leqslant)\cong(b,\subset)$. {\hfill $\qedsymbol$}

\vspace{10pt}

\hypertarget{sec:定理1.3}{}
\noindent\textbf{定理1.3}

非空有限集上的偏序一定有极小(大)元.

\noindent{\kaishu 证明.}

任取偏序结构$(a,\leqslant)$, 假设$a$是非空有限集, 再假设$a$没有极小元.

对任意的$x\in a$, 选取$f(x)\in a$使得$f(x)<x$.
再任取$x_0\in a$, 据此递归定义$p:\omega\to a$ :
\vspace{3pt}
\[\forall n\in\omega:\quad p(n)=\left\{\begin{aligned}
    &x_0,\quad&&n=0,\\[-3pt]
    &f(p(n-1)),\quad&&n>0.
\end{aligned}\right.\\[3pt]\]
则对任意自然数$n$有$p(n+1)=f(p(n))<p(n)$, 显然$p$是单射,
从而$\omega\preccurlyeq a$, 矛盾.

所以$a$有极小元, 同理有极大元. \hfill $\qedsymbol$

\end{document}